
où \(X'\) est le \(X\)-schéma
obtenu en faisant éclater \(Y\). On suppose en outre que \(X\) est quasi-compact et a un faisceau ample.

\textbf{Proposition 4.2.} \uline{Le morphisme} \(f\) \uline{est projectif}, \uline{d'intersection complète}, \uline{et de dimension relative virtuelle nulle}. \uline{On a de plus la relation} :
\[
 f_{\mathrm{gr}}\circ f^{\mathrm{gr}}=\operatorname{Id}_{\Gr K^{\bullet}(X)_{\mathbb Q}} .
\]

La première assertion découle immédiatement de (VII 1.8).

La deuxième relation est conséquence de la relation :
\[
 f_*\circ f^*=\operatorname{Id}_{K^{\bullet}(X)},
\]
montrée dans (VII 3.6).

\textbf{Proposition 4.3.} \uline{Pour tout} \(y\in K^{\bullet}(Y)_{\mathbb Q}\), \uline{on a la relation} :
\[
 f^{\mathrm{gr}}\circ i_{\mathrm{gr}}(y)
 =
 j_{\mathrm{gr}}\bigl(g^{\mathrm{gr}}(y)\cdot c^{d-1}(\check F)\bigr),
\]
\uline{où} \(c^{d-1}\) \uline{désigne la} \((d-1)\)-\uline{ième classe de Chern}.

Cette relation est contenue dans VII 4.8, car elle n'est autre que \(v^{\mathrm{gr}}\circ u^{\mathrm{gr}}=0\).

\textbf{Proposition 4.4.} \uline{Soit} \(X\) \uline{un schéma quasi-compact}. \uline{Si} \(\underline{E}\) \uline{est un faisceau localement libre de type fini sur} \(X\), \uline{de classe} \(E\) \uline{dans} \(K^{\bullet}(X)\), \uline{et de rang} \(n\), \uline{on a} :
\[
 \ch(\lambda_{-1}(E))=c^n(\check E)(\Todd(\check E))^{-1}.
\]

Les deux membres de la relation à démontrer sont multiplicatifs : c'est clair pour \(\ch\circ\lambda_{-1}\), et pour \(\Todd\); c'est aussi vrai pour \(c^n\), qui est la classe de \((-1)^{\varepsilon(E)}\lambda_{-1}(E)\). Par suite, on peut, en passant au fibré en drapeaux de \(\underline{E}\), supposer que \(E\) est somme de classes de faisceaux inversibles, et se ramener à montrer la formule pour un faisceau inversible. Soit
